\newcommand{\DoNotLoadEpstopdf}{}
\documentclass[11pt,reqno]{amsart}
\usepackage[T1]{fontenc}
\usepackage{lmodern,amsmath,amssymb,amsthm,mathtools,mathrsfs,microtype}
\usepackage[margin=1.1in]{geometry}
\usepackage{xr-hyper}
\usepackage[hidelinks]{hyperref}
\externaldocument[supp-]{build/supplement}[supplement.pdf]
\hypersetup{pdftitle={Collision records, arithmetic, and raw local inversion in a triangular Lorentz gas},pdfauthor={Qian Qi}}
\setlength{\emergencystretch}{2em}
\allowdisplaybreaks[2]
\numberwithin{equation}{section}
\newtheorem{theorem}{Theorem}[section]
\newtheorem{lemma}[theorem]{Lemma}
\newtheorem{proposition}[theorem]{Proposition}
\newtheorem{corollary}[theorem]{Corollary}
\theoremstyle{definition}\newtheorem{definition}[theorem]{Definition}
\newtheorem{inputstatement}[theorem]{Input statement}
\theoremstyle{remark}\newtheorem{remark}[theorem]{Remark}
\newcommand{\R}{\mathbb R}\newcommand{\Z}{\mathbb Z}
\newcommand{\one}{\mathbf1}\newcommand{\dd}{\,\mathrm d}
\newcommand{\esssup}{\operatorname*{ess\,sup}}
\newcommand{\diver}{\operatorname{div}}\newcommand{\TV}{\operatorname{TV}}
\title[Collision records and raw local inversion]{Collision records, arithmetic, and raw local inversion in a triangular Lorentz gas}
\author{Qian Qi}
\date{October 13, 2026. A2-DYN, revision 103}
\subjclass[2020]{37D50, 37A50, 60F05, 28A75}
\keywords{Lorentz gas, exact collision records, positive paired traces, endpoint-distance coarea, critical entrances}
\begin{document}
\raggedbottom
\begin{abstract}
We study the exact-label pointwise raw local limit problem for actual
returns in a triangular Lorentz gas. The remaining upper source is
the positive minimum of the two weighted endpoint traces on each
short regular segment. For this same measure we construct two
comparisons preserving distance in roof time from one endpoint,
rather than relative position between the endpoints. Their coarea
densities contain neither a residence factor nor a relative-position
factor. An expanding-endpoint criterion therefore admits all relative
positions, including approach to a positive backward-critical trace
and to a fold of the relative-position map. After physical summation,
the certified source has essential height $O(B^{-1/192})$ under the
original outer schedule. The two clocks also recover part of the
previous failure source with height $O(B^{-1/384})$. The residual
which fails all three comparisons remains positive and is not proved
to vanish. The full pointwise target, exact labels, original guard,
arithmetic classes and preceding proofs are retained.
\end{abstract}
\maketitle
\input{journal-problem-v93}

\subsection{Two additional directions on the positive paired measure}
Write $a(\theta)<u<b(\theta)$ for an original maximal regular
segment, $\ell=b-a$, and
$c^\varepsilon=\min\{f^\varepsilon(a+),f^\varepsilon(b-)\}$.
The earlier transverse comparison fixes $s=(u-a)/\ell$ and has
speed $(1-s)a'+sb'$. At a backward critical entrance with $a'=0$
this speed degenerates as $s\downarrow0$; it may also vanish when
the two endpoint derivatives have opposite signs.

We instead fix either $u-a$ or $b-u$. The exact paired densities
in the two resulting roof coordinates are $c^\varepsilon/|a'|$
and $c^\varepsilon/|b'|$. The absence of $\ell$ and $s$ is an
identity of measures, not an asymptotic cancellation. A one-sided
expansion of $\ell$ keeps every original interior position inside
its segment, even when the starting residence is arbitrarily short.
The comparison still moves between actual trajectories of the same
exact event; it does not continue the radial orbit beyond its first
exit or change any section decision.

Theorem~\ref{thm:v103-gap-height} gives a finite-count pointwise
comparison after physical sheets have been summed. Theorem~\ref{thm:v103-certificate}
and Corollary~\ref{cor:v103-corners} prove its uniformity in relative
position, with the precise reach and trace-variation hypotheses.
Theorem~\ref{thm:v103-recovery} pays the recovered positive source
and retains the exact remainder. In particular a conditional endpoint
certificate is not identified with admission of every physical pair.

The original barycentric comparison is stated first with its proof.
It remains available unchanged and is combined with, not replaced by,
the two endpoint-distance comparisons. The additional critical-trace
calculation distinguishes a positive critical entrance from a genuinely
grazing end. All clocks below use pullback Jacobians for the paired
measure; their integral covariance is printed explicitly.

\input{journal-inputs-v102}
\input{core/244_barycentric_positive_trace_coarea}
\input{core/245_transverse_clock_and_pointwise_trace_height}
\input{core/246_transverse_certificates_and_scalar_schedule}
\input{core/247_endpoint_distance_coarea_and_critical_traces}
\input{core/248_gap_clock_pointwise_recovery}
\input{core/249_expanding_endpoint_certificate}

\appendix
\input{appendices/v103_dependency_map}
\input{journal-references-v100}
\end{document}
